\section{洪荒装置：把路线变成一台机器}

本章看能量奇点的第一台装置。洪荒 70 的价值不是直接追求高 Q，而是验证“用高温超导材料建造完整托卡马克装置”的工程可行性。杨钊用了一个船的类比：过去所有船都用木头造，现在有人说要用钢造船；钢船的浮力原理不变，但焊接、除锈、结构和下水风险都变了。高温超导托卡马克也是这样：聚变物理不变，但主体材料改变后，设计、加工、接口和系统集成全都要重做。

一台装置从想法到运行，要经历物理设计、概念设计、工程设计、制造加工、子系统验收、总装联调和实验运行。洪荒 70 从四个创始人起步，到建成时约百人团队；过程中大量问题不是“公式不会”，而是不知道哪个环节出了问题、哪个参数真实落在什么范围、如何在无法完全测量的情况下做判断。

\begin{figure}[H]
\centering
\includegraphics[width=0.96\textwidth]{fusion-startup-roadmap.png}
\caption{核聚变创业路线图：从物理验证到示范电厂，再到商业电站。自制概念图，依据 01:00:20--01:20:40 对谈内容整理。}
\end{figure}

\begin{knowledgebox}{读图：路线图不是一个大跳跃}
从 Physics 到 Device，再到 Pilot、Grid 和 Commercial，表达的是分阶段降风险。洪荒 70 验证完整高温超导装置工程可行性；洪荒 170 目标是成本受控地实现 Q 大于 10；洪荒 380 才更像完整示范电站和产品。
\end{knowledgebox}

\subsection{第一等离子体是什么}

本节从“路线图”进入“机器第一次活起来”的时刻。洪荒 70 并不是为了直接冲击 Q 大于 10，而是为了证明一台全高温超导托卡马克可以按工程设计完成联调。理解第一等离子体，能避免把它误读成“已经能发电”，也能理解它为什么仍然是极关键的里程碑。

第一等离子体是托卡马克装置真正“点亮”的标志：真空室内的中性气体通过预电离、加热和磁场控制变成等离子体，并呈现出接近设计的磁场位形。它不代表 Q 值很高，也不代表电站已经可用；它代表装置的关键子系统可以联合运行，至少能够把中性气体变成受控等离子体。

\begin{importantbox}{第一等离子体的意义}
第一等离子体不是商业发电指标，而是系统集成指标。它说明磁体、真空、加热、充气、诊断和控制等系统满足基本联动条件，装置从零件集合变成了可运行机器。
\end{importantbox}

\begin{knowledgebox}{工程流程：一台托卡马克怎么从图纸变成机器}
\begin{tabularx}{\textwidth}{p{0.18\textwidth}p{0.34\textwidth}X}
\toprule
阶段 & 主要任务 & 失败风险 \\
\midrule
物理设计 & 确定等离子体目标参数 & 目标过高或约束不清，会导致后续系统失配。 \\
概念设计 & 分解磁体、真空、低温、加热、诊断等系统指标 & 子系统接口冲突会在后期放大。 \\
工程设计 & 出图纸、选型、仿真、工艺验证 & 小样与实物参数不一致。 \\
制造验收 & 加工磁体、真空室、电源、低温系统并逐项验收 & 单件性能衰减或运输/装配损伤。 \\
总装联调 & 多系统一起按工况运行 & 最难定位的是“不知道哪里出了问题”。 \\
\bottomrule
\end{tabularx}
\end{knowledgebox}

\subsection{今日磁体与洪荒 170}

访谈中提到的“今日磁体”是洪荒 170 环向场磁体的重要验证。洪荒 170 需要达到远高于 70 的磁场参数，例如 23 特量级；因此团队先用一个高参数大孔径磁体验证设计、工艺和加工路线。这个磁体不会直接装到 170 上，但它证明：团队有能力制造出未来装置最关键、最贵、最难的子系统。

\begin{knowledgebox}{系统分解：把大风险拆成子系统风险}
\begin{tabularx}{\textwidth}{p{0.20\textwidth}p{0.36\textwidth}X}
\toprule
层级 & 验证什么 & 代表问题 \\
\midrule
洪荒 70 & 高温超导完整装置能否建成并运行 & 材料换代后的系统集成风险。 \\
今日磁体 & 高场大孔径磁体能否制造 & 高磁场、高工程电流密度、结构受力和低温冷却。 \\
洪荒 170 & 低成本装置能否实现 Q 大于 10 & 从工程可行性走向商业可行性的关键台阶。 \\
洪荒 380 & 示范电站和产品化 & 长时间运行、取热发电、客户交付和电站经济性。 \\
\bottomrule
\end{tabularx}
\end{knowledgebox}

\begin{figure}[H]
\centering
\includegraphics[width=0.96\textwidth]{fusion-bottlenecks.png}
\caption{核聚变关键瓶颈：科学、工程、材料、资金和监管同时过线。自制概念图，依据 01:03:27--01:29:59 对谈内容整理。}
\end{figure}

\begin{knowledgebox}{读图：瓶颈不是单点突破}
等离子体、材料、磁体、热工、资金和监管每一项都可能卡住商业化。读图时要注意“同时过线”：某个子系统世界第一不等于电站可用，资金足够也不等于工程风险消失。
\end{knowledgebox}

\subsection{本章小结}

洪荒 70 的意义是把高温超导托卡马克从想法变成完整装置；今日磁体的意义是验证 170 的关键子系统能力；洪荒 170 和 380 则分别对应 Q 大于 10 与示范电站商业化。

